\subsection*{14.3 Prokhorov Theorem and Sequential Compactness}

In the previous two sections, we proved that Condition 1 implies Condition 2, and

Condition 3 implies Condition 4 in Levy's continuity theorem. The reason why Con-

dition 2 implies Condition 3 is characteristic function's continuity (Theorem 9.3)It

remains to prove that Condition 1 follows from Condition 4. The final step requires

the Prokhorov theorem, which we state without proof as follows.

\begin{thmbox}{14.6 (Prokhorov)}
\end{thmbox}

If.(Fn)\infty

n=1 is a sequence of tight cumulative distribution functions, then there

exists a subsequence converging in distribution to the cumulative distribution

of a probability measure.

The Prokhorov theorem is analogous to the Bolzano-Weierstrass theorem. In

real analysis, a subset S in Rn is said to be sequentially compact if any infinite

sequence in S has a convergent subsequence that converges to an element in SThe

Prokhorov theorem says that the space of tight distributions with topology defined

by weak convergence is sequentially compact. For a proof of this theorem, see [ 4,

Thm 3.2.12 and 3.2.13].

We now complete the proof of Levy's theorem. Consider a sequence .(Pn)n\geq1 of

probability measure that is a tight, and let.(Fn)n\geq1 be the corresponding cumulative

distribution functions. In Levy's continuity theorem, the characteristic function

\phin(t) of Pn is assumed to converge pointwise, with the limit function denoted by

\phi(n).

By the Prokhorov theorem, there exists a subsequence .(an)n\geq1 of N such that

Fan(x) converges pointwise to the cumulative distribution function F(x) of a

probability measure P at every continuity point of F(x) Using Theorem 14.3

and the equivalence of weak convergence and convergence in distribution, we can

show that the characteristic functions\phian (t) converge pointwise to the characteristic

function of P. Since we assume that limn\to\infty \phian(t)= \phi(t) , we can conclude that

\phi(t) is the characteristic function of P.

We now demonstrate that the entire sequence .(Pn)n\geq1 of probability measures

weakly converges to PLeth(x) be a bounded and continuous function. We must

show that

limn\to\infty

\int

R

h(x) dPn(x)=

\int

R

f (x) dP(x).

To prove this, we use the following fact from real analysis:

Given a numerical sequence .(\alphan)n\geq1, if every subsequence of .(\alphan)n\geq1 contains

a further subsequence that is convergent to \beta, then the original sequence .(\alphan)n\geq1 is

convergent to \beta.

Let i1, i2,i 3,.be a subsequence of N. We consider the subsequence

( \int

h(x) dPij (x)

)

j\geq 1

(14.8)

Since any subsequence of a tight sequence is tight, the sequence of probability

measures .(Pij )j\geq 1 is tight. By Prokhorov theorem, there exists a sub-subsequence

ijk , where j1 <j 2 <j 3 <\cdot\cdot\cdot , such that Fijk (x) converges pointwise to

the distribution function G(x) of a probability measure Q at every continuity

point of G(x)By Theorem 14.3, the characteristic functions associated to this

sub-subsequence converge to the characteristic function of Q , which is also the

characteristic function of P due to the assumption that limk\to\infty \phiijk (t)= \phi(t) .

By the inversion formula in Theorem 9.5, Q and P have the same distribution.

Hence,

( \int

h(x) dPijk (x)

)

k\geq 1

is a sub-subsequence of (14.8) converging to

lim

k\to\infty

\int

h(x) dPijk (x)=

\int

h(x) dQ(x)=

\int

h(x) dP(x).

Therefore,.

\int

hdP n converges to.

\int

hdP asn\to\infty , and this completes the proof of

Levy's continuity theorem.

We illustrate Levy's continuity theorem with two classic examples. The first

example involves discrete random variables converging to another discrete random

variable, while the second example involves discrete random variables converging

to a continuous random variable.

\textbf{Example 14.3.1 (Binomial Converging to Poisson)} \\

Let Xn be a random variable with binomial distribution Bin(n, pn) with pn =\lambda/n ,f o rafi x e d

positive constant. \lambda. The characteristic function of. Xn is

\phiXn (t)= ( 1- p n +p neit )n.

If we taken\to\infty ,

limn\to\infty \phiXn (t)= limn\to\infty

(

1- \lambda

n(1- e it )

) n

=e \lambda(eit -1),

which is the characteristic function of a Poisson random variable. By Theorem 14.1,Xn converges

to a Poisson random variable with mean. \lambda in distribution.

\textbf{Example 14.3.2 (Geometric Converging to Exponential)} \\

Let \lambda be a positive constant and Xn be geometrically distributed with success probability pn =

\lambda/n,f o rn= 1 ,2,3,.The characteristic function of. Xn is

\phiXn (t)= pneit

(1- ( 1- p n)eit ).

LetYn =X n/n. We can check that\phiYn (t)= \phi Xn (t/n) converges pointwise asn\to\infty ,

limn\to\infty \phiYn (t)= limn\to\infty

eit/n \lambda/n

1- ( 1- \lambda/n)e it/n = \lambda

\lambda- it .

By Theorem 14.1, we have weak convergenceYn

D

-\to Exp (\lambda).
